\section{Results and Discussion}
\label{sec:results}

The methodology was applied end-to-end to the conceptual design of a
small electric tail-sitter VTOL UAV inspired by early V-BAT-class
demonstrators: a single ducted fan in the aft fuselage, four jet vanes
in the fan efflux providing pitch/yaw/roll authority in hover and
transition, wing ailerons as the cruise-phase roll effector, and
tail-down vertical takeoff and landing. The mission requirement is a
0.5\,kg payload carried through a deliberately short-hover profile ---
120\,s of vertical flight, 40\,s of transitions (billed at hover
power), and 900\,s of cruise at 20\,m/s ($\approx$18\,km range) with a
20\% energy reserve. This section reports (i) the converged design
outcome, (ii) the design-point evolution that produced it, (iii)
process metrics extracted from the project's git history quantifying
the AI's role, and (iv) discussion against the research question.

\subsection{Converged design point}
\label{sec:design-results}

Closure is performed with transparent first-order physics: hover power
from momentum (actuator-disk) theory with an overall figure-of-merit
efficiency chain, cruise power from a conservative lift-to-drag model,
and battery mass from mission energy through pack-level specific
energy, with the pack's Joule losses folded in self-consistently (the
battery efficiency is the mission-averaged $I^2R$ efficiency of the
nominal 90\,m$\Omega$ pack, iterated to the closure fixed point).
Table~\ref{tab:designpoint} summarizes the converged vehicle;
Fig.~\ref{fig:vehicle} shows the parametric solid model regenerated by
the pipeline on every run.

\begin{table}[htbp]
\centering
\caption{Converged conceptual design point (v0.5.2 snapshot).}
\label{tab:designpoint}
\footnotesize
\setlength{\tabcolsep}{4pt}
\begin{tabular}{@{}llll@{}}
\toprule
\textbf{Quantity} & \textbf{Value} & \textbf{Quantity} & \textbf{Value} \\
\midrule
MTOW (closure) & 2.518\,kg & Wing loading & 131.2\,N/m$^2$ \\
All-up mass (as-selected) & 2.336\,kg & Wing area / span & 0.1883\,m$^2$ / 1.063\,m \\
Payload & 0.5\,kg & Aspect ratio & 6.0 \\
Hover electrical power & 743\,W & Airfoil & NACA 4412 ($C_{L,\max,3D}=1.489$) \\
Peak battery discharge & 8.1\,C & Cruise $C_L$ / $L/D$ & 0.535 / 13.35 \\
Rotor & 203\,mm, 3-blade $8\times6$ & Cruise speed & 20\,m/s ($\approx$1.02\,$V_{\mathrm{MD}}$) \\
Blade section / solidity & Clark~Y / $\sigma=0.077$ & Mission & 120\,s hover + 40\,s trans.\ + 900\,s cruise \\
1/rev forcing & $\sim$204\,Hz & Battery (frozen) & Li-ion 6S1P, 108\,Wh, 450\,g \\
Fuselage envelope & $\varnothing$107 $\times$ 533\,mm & Structure & segmented-FDM semi-monocoque \\
\bottomrule
\end{tabular}
\end{table}

\begin{figure}[htbp]
\centering
\begin{subfigure}[t]{0.48\textwidth}
  \centering
  \includegraphics[width=\linewidth]{vbat_render}
  \caption{Rendered vehicle, tail-down on its landing legs.}
\end{subfigure}\hfill
\begin{subfigure}[t]{0.48\textwidth}
  \centering
  \includegraphics[width=\linewidth]{fuselage_layout}
  \caption{Fuselage packaging and bay stack from the layout stage.}
\end{subfigure}
\caption{Pipeline-generated design artifacts: every figure is
regenerated from configuration by the deterministic pipeline; none is
hand-drawn.}
\label{fig:vehicle}
\end{figure}

The commercial-component freeze (pipeline stage 12) selected, against
requirements \emph{derived} from the converged design point, a
Pixhawk-class PX4 flight controller, a telemetry-capable 80\,A-class
ESC, the pinned 3-blade $8\times6$ propeller, sub-10\,g-class
control-surface servos, and a 21700 Li-ion 6S pack that undercuts the
sized LiPo pack by $\sim$203\,g --- flipping the as-selected all-up
mass 182\,g \emph{below} the conceptual closure. The as-selected
re-solve (stages 13--15) then re-solved the affected subsystems with
real hardware masses and envelopes and confirmed physical fit in the
unchanged $\varnothing$107\,mm hull.

Crucially, the pipeline also reports what does \emph{not} close.
Standing findings pinned at the final snapshot include: a marginal ESC
cold plate (its $\sim$37\,W hover load needs a plate heavier than the
ESC mass allocation, with only a few $^{\circ}$C margin); a battery
pack that, at the 40\,$^{\circ}$C design ambient and nominal
90\,m$\Omega$ resistance, ends the mission $\sim$5\,$^{\circ}$C over
its 60\,$^{\circ}$C limit (making measured pack resistance a
procurement acceptance criterion); and avionics-bay and motor mass
allocations exceeded by $\sim$25\,g and $\sim$79\,g respectively.
These are carried openly as accepted risks with named retirement
paths, not tuned away --- the behavior the standing-findings
discipline of Section~\ref{sec:verification} was designed to produce.

\subsection{Design-point evolution}
\label{sec:evolution}

Table~\ref{tab:evolution} traces the design point across the tagged
snapshots. The trajectory is deliberately non-monotonic: the design
got \emph{heavier} whenever a model was made more honest (the
3D-printability penalty at v0.2.0; the battery internal-resistance
accounting at v0.5.1, $+215$\,g) and lighter whenever architecture or
procurement evidence justified it (the externally reviewed
semi-monocoque re-baseline at v0.4.0; the larger, lower-disk-loading
propeller the component freeze forced at v0.4.1). Each row is a
complete, internally consistent, released design snapshot with its
regression pins, CFD references, CAD exports, and changelog updated in
the same change --- there is no intermediate state in the history in
which the artifacts disagree.

\begin{table}[htbp]
\centering
\caption{Design-point evolution across tagged snapshots. Every row is
a fully re-converged, regression-pinned release.}
\label{tab:evolution}
\small
\begin{tabularx}{\textwidth}{@{}llllX@{}}
\toprule
\textbf{Tag} & \textbf{Date} & \textbf{MTOW} & \textbf{Hover} &
\textbf{Driver of the change} \\
\midrule
v0.1.0 & Jul 09 & 2.50\,kg & $\sim$770\,W &
Initial closure: NACA 2412 wing, 195\,mm COTS fan assumption \\
v0.2.0 & Jul 10 & 3.06\,kg & $\sim$1030\,W &
Honest 3D-printed-structure penalty enters the closure; vibration
isolation and modularity split lines added \\
v0.3.0 & Jul 10 & 3.06\,kg & --- &
Thermal paths sized as structure; marginal ESC cold plate first
reported as a standing finding \\
v0.4.0 & Jul 11 & 2.376\,kg & 710\,W &
Semi-monocoque clamshell fuselage from independent expert review;
structural budget re-baselined to an explicit member model \\
v0.4.1 & Jul 15 & 2.302\,kg & 652\,W &
Component freeze reveals the 195\,mm fan class is unbuyable at this
scale; propulsion decision amended by ADR to a 203\,mm 3-blade
propeller in the airframe duct \\
v0.4.5 & Jul 17 & 2.303\,kg & --- &
Li-ion 21700 pack wins the battery freeze; battery mission-transient
thermal check added from external review \\
v0.5.0 & Jul 17 & 2.303\,kg & --- &
Stall-speed inconsistency found by expert review (placeholder
$C_{L,\max}$ survived in one report); wing re-sized to the real
airfoil limit \\
v0.5.1 & Jul 19 & 2.518\,kg & 743\,W &
NACA 4412 re-selection exploits the now-honest stall constraint;
pack $I^2R$ losses folded into closure ($+215$\,g accepted) \\
v0.5.2 & Jul 21 & 2.518\,kg & 743\,W &
Versioned aerodynamic hand-off contract iterated with the external
solver consumer (six schema releases) \\
\bottomrule
\end{tabularx}
\end{table}

\subsection{Process evidence: quantifying the AI's role}
\label{sec:process-results}

Because every AI-implemented commit carries a machine-readable
co-authorship trailer, the division of labor is recoverable directly
from the git history, providing unusually direct process evidence.

\paragraph{Two phases, one repository.} The project history contains a
natural comparison. In a first, entirely manual phase (4~February --
27~March), 28 commits over seven calendar weeks carried the study
through mission sizing, initial take-off-weight iteration, airfoil
selection, and control-vane design --- roughly the first three of the
eventual sixteen pipeline stages, with physics gradually migrating out
of notebooks. In the AI-assisted phase (4--24~July), 109 commits over
three calendar weeks of part-time work delivered the remaining
thirteen pipeline stages plus the full governance and verification
infrastructure: parametric CAD, CFD cases, vibration, thermal, mass
properties, wiring, component freeze, the as-selected re-solve,
sixteen executed notebooks in CI, 147 pinned regression tests, and
thirteen tagged design releases (v0.1.0 through v0.5.2 in twelve
days). While the two phases are not perfectly controlled (the second
benefited from the first's groundwork), the contrast in scope per
calendar week is roughly an order of magnitude and is consistent with
the acceleration claims in the surrogate-model literature, achieved
here at the workflow level rather than inside a single discipline.

\paragraph{Contribution split.} Of the 109 AI-phase commits, 70
(64\%) carry AI co-authorship trailers; the remainder are human-only
commits (merges, reviews, manual fixes, procurement research). All 23
engineering pull requests of the AI phase were human-reviewed before
merge. The final repository comprises $\sim$11{,}200 lines of
deterministic physics library code, $\sim$1{,}800 lines of tests, 18
commented configuration files, 16 orchestration notebooks, and 18
architecture decision records --- the ADRs and changelog entries
drafted by the agent as part of the changes they document.

\paragraph{Error-catching channels.} The history also documents
\emph{which} safety net caught each class of error, summarized in
Table~\ref{tab:errors}. The pattern is consistent with the layered
design of Section~\ref{sec:method}: automated pins caught silent
drift, schema-governed hand-offs caught cross-tool integration faults,
and independent human expertise caught conceptual errors invisible to
any automated check --- confirming that expert review is a
load-bearing element of the methodology, not a formality.

\begin{table}[htbp]
\centering
\caption{Representative faults during the case study and the
methodology layer that caught each.}
\label{tab:errors}
\small
\begin{tabularx}{\textwidth}{@{}XX@{}}
\toprule
\textbf{Fault} & \textbf{Caught by} \\
\midrule
Propulsion/vane CFD cases silently meshing an obsolete design point &
Cross-consumer consistency discipline; cases re-sourced from hand-offs
with a design-point stamp that refuses stale reuse \\
\addlinespace
Stall requirement silently unmet (placeholder $C_{L,\max}$ vs.\
selected airfoil) & Independent expert review (v0.5.0) \\
\addlinespace
Unrealistically optimistic battery efficiency (implied
$\sim$23\,m$\Omega$ pack) & Expert review prompting the transient
model; closure re-derived from pack physics (v0.4.5--v0.5.1) \\
\addlinespace
Wing control points inside the fuselage in the external solver &
Versioned hand-off schema iteration with the consumer (placement
anchor made mandatory from schema 1.5.0) \\
\addlinespace
Unbuyable propulsion assumption (195\,mm fan class) & Deterministic
component-freeze stage; assumption amended by ADR (v0.4.1) \\
\addlinespace
Refactoring risk while moving physics out of notebooks &
Byte-identical hand-off verification (ADR-0013, ADR-0018) \\
\bottomrule
\end{tabularx}
\end{table}

\subsection{Discussion}
\label{sec:discussion}

\paragraph{Answering the research question.} The case study supports
an affirmative, but conditional, answer to the RQ. The AI accelerated
the program by roughly an order of magnitude in delivered scope per
calendar week --- but the acceleration was \emph{not} obtained by
letting the model produce engineering answers. Every number in
Table~\ref{tab:designpoint} traces to deterministic, configured,
regression-tested physics; the AI's leverage came from (i)
implementing multidisciplinary changes completely, so that a design
move lands with its pins, consumers, decision record, and changelog in
one reviewable diff; (ii) eliminating the bookkeeping latency that
normally separates a design decision from a fully consistent design
state; and (iii) making honest-model refinements cheap enough that
unfavorable results (a $+215$\,g mass growth, a marginal cold plate)
were adopted rather than resisted.

\paragraph{Division of labor.} The observed equilibrium assigns the
AI what LLM agents are reliably good at --- exhaustive propagation,
code implementation, documentation, consistency --- and reserves for
humans what they were uniquely necessary for in this program:
requirement setting, acceptance of every decision, procurement
judgment, and conceptual review. It is notable that both of the case
study's most consequential errors (the stall inconsistency and the
optimistic battery efficiency) were latent \emph{before} the AI phase
and were surfaced during it, one by human expert review and one by an
expert-prompted model refinement the AI then implemented. Neither the
AI alone nor the automation alone would have sufficed; the methodology
is precisely the structure that lets each catch what the other misses.

\paragraph{Generality.} Nothing in
Sections~\ref{sec:pipeline}--\ref{sec:iteration} is specific to
tail-sitters: the constitution, one-way pipeline, pins, ADR
discipline, and interaction protocol transfer directly to any
configuration-driven engineering program with computable analysis
stages. The tail-sitter case exercised the framework hard ---
hover/cruise coupling, dual control effectors, thermal and vibration
constraints, and a procurement reality check all interact --- which is
exactly why it was chosen as the validation vehicle.

\paragraph{Limitations.} The study is a single case with a team
already expert in the domain and the tooling; the phase comparison is
suggestive, not controlled. All analyses are conceptual-level
(momentum theory, lifting-line-class aerodynamics, first-order thermal
and structural models); transition dynamics are energetically budgeted
but not simulated, and no hardware has flown --- the PX4-based flight
validation and the external VLM/BEMT and CFD loops are the designed
next steps, wired into the same hand-off contracts. Finally, the
methodology's guardrails bound, but cannot eliminate, LLM
fallibility: the protection is structural (nothing merges without
passing pins, CI, and human review), and the residual risk
concentrates exactly where automated checks are weakest ---
conceptual assumptions --- which is why independent expert review
remains mandatory in the loop.
